% !TeX root = ../../Book3.tex
\chapter{Weierstrass 乘积、Mittag–Leffler 与 Hadamard 分解}
\label{b3:ch:20}
\BookChapterNumber{b3:ch:20}
% 本章摘要（不计入编号）
\begin{chapterabstract}
无穷乘积通过扣除对数的低阶项实现指定零点，Mittag--Leffler 加法构造则实现指定主要部分。两种构造分别留下无零点乘法因子与整函数加法项。有限阶增长进一步将乘法因子压成指数多项式，得到 Hadamard 分解及正弦的典范乘积。
\end{chapterabstract}

% 本章目标（不计入编号）
\begin{learninggoals}
能检验正规乘积和零点阶数；能按离散序列选择初等因子；能构造整平面亚纯函数及整插值；能说明 Hadamard 证明中选圆周、乘积下界和实部控制的作用。
\end{learninggoals}

如果要在正整数处各放一个零点，直接相乘 \(1-z/n\) 会留下发散的一次对数项。把它改成 \((1-z/n)e^{z/n}\)，二次尾项就可以求和。一个看似小的指数修正使零点构造有了稳定的极限；而有限阶函数又反过来限制这些修正能增长多少。

\ProseParagraphBreak

本章使用\chapref{b3:ch:19}的零点计数与 Cauchy 估计。进一步阅读为 Conway 的《Functions of One Complex Variable I》\cite{Conway1978Functions}和 Levin 的《Lectures on Entire Functions》\cite{Levin1996Lectures}，可据零点构造与有限阶分解的主题继续阅读。

% !TeX root = ../../Book3.tex
\section{无穷乘积}
\label{b3:sec:2001}

有限个零点可以用多项式乘积实现。无穷个零点要求控制乘积尾部，且必须确认极限除了指定因子之外不会产生额外零点。

\subsection{收敛判据}
\label{b3:subsec:200101}

若 \(\sum_n|u_n|<\infty\)，则 \(\prod_n(1+u_n)\) 收敛；当所有因子非零时极限也非零。证明取尾部 \(|u_n|\le1/2\)，使用 \(\log(1+u_n)=\sum_{k\ge1}(-1)^{k+1}u_n^k/k\) 的固定小圆分支，得到 \(|\log(1+u_n)|\le2|u_n|\)。尾部乘积是收敛对数和的指数，故非零；有限头部另行相乘。

不控制对数和，仅有 \(u_n\to0\) 不足以保证收敛：\(\prod(1+1/n)\) 发散。也不能把“收敛”一律解释为非零；\(\prod(1-1/n)\) 从第二项起趋零。

\subsection{绝对收敛}
\label{b3:subsec:200102}

绝对收敛的乘积允许重排，其原因是小因子的对数绝对可和。无穷乘积的“绝对收敛”通常指 \(\sum|u_n|<\infty\)，不是对乘积模再作一次级数。

若有某个因子为零，整乘积为零，此时对数表述只用于去掉该因子后的尾部。后面在函数论中，零点附近总会把产生零点的有限因子与无零点的尾部拆开。

\subsection{正规收敛}
\label{b3:subsec:200103}

\begin{proposition}\label{b3:prop:20-normal-product}
若全纯函数 \(u_n\) 在区域 \(D\) 上满足对每个紧集 \(K\Subset D\)，\(\sum_n\sup_K|u_n|<\infty\)，则 \(\prod(1+u_n)\) 局部一致收敛到全纯函数；在无因子零点的地方，极限非零。
\end{proposition}
\begin{proof}
固定紧集附近稍大的紧集，尾部因子一致落在 \(|u|\le1/2\)，对数级数由上述估计正规收敛，其和全纯，指数给全纯非零尾部。乘上有限头部，得到局部一致极限。不同紧集的尾部表示来自同一有限乘积序列，故彼此一致。
\end{proof}
局部一致性保证可逐项微分对数：无零点的紧子集上 \(P'/P=\sum u_n'/(1+u_n)\)，导数收敛由全纯局部一致收敛定理给出。

\subsection{零点与乘积}
\label{b3:subsec:200104}

在固定点附近，若只有有限个因子为零，就把它们提出；余下尾部由正规对数和给非零函数。因此极限的零点阶数恰为有限零因子的阶数和，没有额外零点。

如果零因子在内部无限聚集，极限可能恒零，而指定离散零点的任务本身就与恒等定理冲突。本章的零点序列始终要求无有限聚点。


对乘积尾部还可给出误差界。若 \(S_N(K)=\sum_{n>N}\sup_K|u_n|\)，则由有限乘积展开
\[
\left|\prod_{n=N+1}^{M}(1+u_n)-1\right|\le
\prod_{n=N+1}^{M}(1+|u_n|)-1\le e^{S_N(K)}-1.
\]
令 \(M\to\infty\)，这个界仍成立。它直接控制尾部相对误差，而不仅是存在一个极限。有限头部在紧集的最大模有界，所以全乘积的绝对误差亦受同一尾和控制。若需要数值计算，先选 \(N\) 使尾和小，再单独计算头部；没有尾部界的有限截断不能证明函数恒等式。

\ProseParagraphBreak

以 \(u_n(z)=z/n^2\) 为例，\(|z|\le R\) 时 \(S_N\le R/N\)。只要没有碰到有限头部的零点，截断相对误差不超过 \(e^{R/N}-1\)。若零点是当前计算目标，必须提出相应因子后再谈相对误差，因为在零值处除以极限没有意义。相同估计还说明局部一致性不能改成全平面一致性：随着 \(R\) 增大，所需 \(N\) 也增大。

\ProseParagraphBreak

条件收敛的乘积可能依赖顺序。例如实数列 \(u_n=(-1)^n/n\) 从 \(n=2\) 起满足 \(\sum u_n\) 条件收敛、\(\sum u_n^2<\infty\)。由
\(\log(1+u_n)=u_n+O(u_n^2)\)，原顺序乘积收敛为正数，但正负项重排能改变一次和，从而改变乘积。这里所有因子甚至都正；失败来自绝对收敛不足，而不是复杂的复数分支。典范因子就是有目的地消除这些可能依赖顺序的低阶和。


\ProseParagraphBreak

正规收敛还允许检验导数，但须区分乘积与对数的适用位置。若紧集避开所有零点，取稍大的无零点邻域，并在那里对尾部统一选择对数，则 Cauchy 公式给对数级数各阶导数的局部一致收敛。若紧集包含一个零点，先分离该点的有限零因子，然后只对剩下的非零部分取对数。若直接写 \(P'/P=\sum u_n'/(1+u_n)\)，左右都是亚纯函数；在零点处谈全纯一致收敛并无意义。

\ProseParagraphBreak

例如 \(P(z)=\prod_{n\ge1}(1+z/n^2)\) 在原点附近的对数为
\[
\log P(z)=\sum_{k\ge1}\frac{(-1)^{k+1}z^k}{k}\sum_{n\ge1}n^{-2k},\qquad |z|<1.
\]
先取 \(|z|\le r<1\)，双重级数的绝对值由 \(\sum_n\sum_k r^k/n^{2k}\) 控制，后者不超过 \(r(1-r)^{-1}\sum_n n^{-2}\)。于是取原点一次导数得到 \(P'(0)=\sum n^{-2}\)，二次导数则为 \((\sum n^{-2})^2-\sum n^{-4}\)。对数系数记录单因子的贡献，函数系数还包含因子之间的交叉项；把这两类系数混在一起会漏掉平方项。

\ProseParagraphBreak

无穷乘积定理不能仅由各因子趋于一替代。常数因子 \(1+1/n\) 虽趋于一，其前 \(N\) 项的乘积为 \(N+1\)，仍发散。常数因子 \(1-1/n\) 从第二项开始乘积为 \(1/N\)，极限恒为零，不能用来断言“除有限个零因子以外尾部非零”。这两个计算说明，可和的对数尾项同时承担有限性与非消失性两项责任。判别时应检查它控制哪一项，不能只检查逐因子的极限。

% !TeX root = ../../Book3.tex
\section{Weierstrass 典范因子}
\label{b3:sec:2002}

直接相乘 \(1-z/a_n\) 往往失败，因为它的对数从一次项开始。乘上一个指数，把若干低阶项扣掉，尾部就可以变成更高次的可和项。

\subsection{初等因子}
\label{b3:subsec:200201}

定义\term[weierstrass-chudeng-yinzi]{Weierstrass 初等因子}[Weierstrass elementary factor]
\[
E_0(w)=1-w,\quad E_p(w)=(1-w)\exp\left(w+\frac{w^2}2+\cdots+\frac{w^p}p\right),\ p\ge1.
\]
它只有一个单零点 \(w=1\)，且 \(E_p(0)=1\)。指数部分不添加零点。

\subsection{典范因子}
\label{b3:subsec:200202}

当 \(|w|<1\)，固定对数分支给
\[
\log E_p(w)=-\sum_{k=p+1}^{\infty}\frac{w^k}k.
\]
所以 \(|w|\le1/2\) 时 \(|\log E_p(w)|\le2|w|^{p+1}\)，并可由 \(|e^v-1|\le e^{|v|}|v|\) 控制因子偏离一。若 \(\sum|a_n|^{-p-1}<\infty\)，则 \(\prod E_p(z/a_n)\) 在每个紧集正规收敛。这里 \(p\) 固定，而紧集半径只进入有限常数。

\subsection{属}
\label{b3:subsec:200203}

能使上述幂和收敛的最小非负整数 \(p\)，称该典范乘积的\term[dianfan-chengji-shu]{属}[genus]。当收敛指数 \(\lambda\) 不是整数时，\(p=\lfloor\lambda\rfloor\)；整数指数时，还要检查临界幂和是否收敛，不能只根据 \(\lambda\) 的值决定。

例如 \(a_n=n\) 的指数一，\(\sum1/n\) 发散而 \(\sum1/n^2\) 收敛，须用属一因子。若把正负整数成对相乘，一次指数正负抵消，可写成 \(\prod(1-z^2/n^2)\)；这依赖明确的配对，而不是许可随意重排一个条件收敛乘积。


初等因子的指数修正还控制对数导数。在避开 \(a_n\) 的地方，直接微分得到
\[
\frac{\mathrm{d}}{\mathrm{d}z}\log E_p(z/a_n)
=-\frac1{a_n-z}+\sum_{k=1}^{p}\frac{z^{k-1}}{a_n^k}
=-\frac{z^p}{a_n^p(a_n-z)}.
\]
当 \(|z|\le R\)、\(|a_n|>2R\) 时，最后一式的模不超过 \(2R^p/|a_n|^{p+1}\)。所以同一幂和不仅保证乘积存在，也保证尾部对数导数正规收敛。零点附近的留数由有限头部贡献，值为一；这将乘法的零点数据重新转成加法的主要部分数据。

\ProseParagraphBreak

对 \(a_n=n\)，属一乘积
\[
P(z)=\prod_{n\ge1}(1-z/n)e^{z/n}
\]
在原点满足 \(P(0)=1\)、\(P'(0)=0\)，且对数展开为
\[
\log P(z)=-\sum_{k\ge2}\frac{z^k}{k}\sum_{n\ge1}n^{-k},\quad |z|<1.
\]
换序由 \(\sum_{n,k}|z|^k/(kn^k)<\infty\) 保证。尤其二阶系数取决于 \(\sum n^{-2}\)，一次系数却因修正消失。这是初等因子“扣除低阶项”的具体可计算效果。

\ProseParagraphBreak

属在临界指数处的差别也能构造出来。点列 \(a_n=n\log^2(n+1)\) 的计数指数仍为一，但 \(\sum1/|a_n|\) 收敛，所以属零已够。点列 \(a_n=n\log(n+1)\) 的指数同样为一，而一次幂和发散，须用属一。二者的收敛指数相同，最小可用属不同。积分判别中出现的对数因子因此不能在临界问题中随意忽略。


\ProseParagraphBreak

若某个 \(a\) 应有重数 \(m\)，可以将它在点列中重复 \(m\) 次，也可以直接取因子 \(E_p(z/a)^m\)。对正规收敛的检验，后一写法须带上权重 \(m\)：充分条件成为 \(\sum m_n|a_n|^{-p-1}<\infty\)。遗漏重数会改变计数函数和收敛指数。例如在每个正整数 \(n\) 放重数 \(n\) 的零点时，带重数的幂和是 \(\sum n^{1-s}\)，收敛指数为二，属零或属一均不足；属二给 \(\sum n^{-2}\) 的尾和。

\ProseParagraphBreak

关于最小属还有一个容易混淆的区分。绝对收敛判据针对逐个初等因子的乘积；把具有对称性的因子成组后，低次项可能彼此抵消。例如正负整数的属一因子满足
\[
E_1(z/n)E_1(-z/n)=1-z^2/n^2.
\]
因此成对乘积有一个看似属零的二次因子表达，但它并没有使 \(\sum_{n\ne0}|n|^{-1}\) 收敛。解释“属”时必须说清楚使用单个零点的标准初等因子，还是已经经过配对的表达。后一表达式有利于计算，却不能反推前一幂和判据失效。

% !TeX root = ../../Book3.tex
\section{Weierstrass 乘积定理}
\label{b3:sec:2003}

离散性是指定零点的唯一存在性障碍。即使没有任何有限幂和收敛，逐项增加初等因子的阶数仍可使尾部收敛。

\subsection{指定零点}
\label{b3:subsec:200301}

\begin{theorem}[Weierstrass：整平面版本]\label{b3:thm:20-weierstrass}
任意无有限聚点的复平面点列与指定正整数重数，可以实现为一个非恒零整函数的全部零点及重数。
\end{theorem}
\begin{proof}
把非零点按重数重复记为 \(a_n\)，则 \(|a_n|\to\infty\)。取
\(P(z)=z^m\prod_{n\ge1}E_n(z/a_n)\)。对固定 \(|z|\le R\)，充分大 \(n\) 时 \(|z/a_n|\le1/2\)，故对数尾项不超过 \(2\cdot2^{-n-1}\)，一致可和。正规乘积命题给整函数；有限头部与非零尾部分离，保证零点及重数正好指定。
\end{proof}
所得函数可能增长非常快，定理没有给阶或型的最优控制。

\subsection{典范乘积}
\label{b3:subsec:200302}

若零点的收敛指数有限，可取固定 \(p\) 的典范乘积，而不用变动的 \(E_n\)。这种选择减少指数修正并保留增长信息。

给定同一零点列，所有实现函数都形如 \(P e^g\)，其中 \(g\) 整。因为商在每个零点处可去且无零点，商的对数导数有整平面原函数，指数还原该商。加法常数由一个点的商值决定。

\subsection{零点序列}
\label{b3:subsec:200303}

一个点列若在有限点聚集，任何包含它的整函数都由恒等定理恒零；这说明“无有限聚点”既必要又充分。有限列给多项式，无穷列给真正的超越整函数。

指定重数并不指定导数系数。若还要在离散点给出函数值或有限阶导数，需要插值；下一节用主要部分给出这种额外自由度的构造。


若指定的零点极为缓慢地趋于无穷，例如 \(a_n=\log(n+1)\)，任何固定 \(p\) 的幂和都发散。Weierstrass 定理仍适用：对紧集 \(|z|\le R\)，从足够大的 \(n\) 起有 \(|z/a_n|\le1/2\)，变动阶数 \(p_n=n\) 使误差指数衰减。这个例子区分了“离散性保证存在”与“有限收敛指数保证有限属”。对应零点计数约为 \(e^r\)，所以任何实现它们的非恒零整函数都必须是无穷阶；这由上一章的零点密度限制直接得到。

\ProseParagraphBreak

相反，零点 \(a_n=n^2\) 的属零乘积具有较慢增长。考虑直接由幂级数定义的整函数
\[
F(z)=\sum_{k\ge0}\frac{(-1)^k\pi^{2k}z^k}{(2k+1)!}
=\frac{\sin(\pi\sqrt z)}{\pi\sqrt z}
\]
的后一表达式按前面级数解释，不需要在全平面选择平方根。更换平方根的符号不改变商，所以局部表达也相容。它的零点恰为正整数平方，每个为单零点；在负实轴 \(z=-r\) 上由 \(\sinh(\pi\sqrt r)\) 得增长下界，复平面上由指数表示得同阶上界，故阶为二分之一。一个非整数阶整函数因而不必靠抽象存在定理展示，可以直接从偶幂级数得到。

\ProseParagraphBreak

对于有限零点集，分解 \(f=P e^g\) 中 \(P\) 是普通多项式。若 \(f\) 又有多项式增长，\(g\) 的自由度实际上只能给常数因子；若只要求有限阶，则 \(g\) 可为多项式。若没有任何增长条件，\(g=e^z\) 等无穷阶例子仍可出现。零点集合本身无法区分这些函数，必须把无穷远条件作为独立数据。


\ProseParagraphBreak

构造的局部性可以用圆盘上的有限头部说明。给定 \(R\)，离散性使 \(|a_n|\le2R\) 的零点只有有限多个。它们的因子可能很大，但全是已固定的整函数；收敛只由其余尾部负责。对 \(|a_n|>2R\)，取 \(p_n=n\) 后
\[
\sup_{|z|\le R}|\log E_n(z/a_n)|\le2(R/|a_n|)^{n+1}\le2^{-n}.
\]
这给出了无需根据每个紧集重新排序的选择。所有半径使用同一个乘积，只是各自的有限头部长度不同。零点集合若在平面内有聚点，上述有限性马上失去；恒等定理同时说明非恒零整函数也不可能实现这类数据。

\ProseParagraphBreak

存在性没有给出规范化的唯一性。指定 \(P(0)=1\) 只消除了常数乘法自由度，仍可乘上 \(e^{az}\)；再指定 \(P'(0)=0\) 也仍可乘上 \(e^{bz^2}\)。只有增长约束或更多结构条件才会继续减少自由度。在有限阶一的情形，两个同零点函数分别有指数一次多项式的分解，因此它们的商为 \(e^{az+b}\)；这时指定原点值及一次对数导数才足够。若原点是零点，应先除去共同的原点零因子，再施加这些条件。

% !TeX root = ../../Book3.tex
\section{Mittag–Leffler 定理}
\label{b3:sec:2004}

乘法构造指定零点，加法构造则指定极点附近的负幂项。把每个远处极点的低阶 Taylor 项扣掉，可以沿用同样的尾部控制思想。

\subsection{主要部分}
\label{b3:subsec:200401}

亚纯函数在极点 \(a\) 的 Laurent 负幂有限和 \(Q_a(z)=\sum_{k=1}^{m_a}c_{a,k}(z-a)^{-k}\) 称为\term[zhuyoubufen]{主要部分}[principal part]。剩余部分在 \(a\) 全纯。指定主要部分比只指定极点阶数强，因为它固定了每个负幂系数及留数。

\subsection{亚纯插值}
\label{b3:subsec:200402}

设 \(a_n\to\infty\)，给定主要部分 \(Q_n\)。选 \(R_n<|a_n|\)，且 \(R_n\to\infty\)，在圆盘 \(|z|\le R_n\) 上用 \(Q_n\) 关于零的 Taylor 多项式 \(T_n\) 逼近到误差 \(2^{-n}\)。这可行，因为 \(Q_n\) 在半径 \(|a_n|\) 内全纯，Taylor 级数在更小闭圆盘一致收敛。

于是 \(\sum(Q_n-T_n)\) 在避开极点的每个紧集上正规收敛。每个指定极点附近，除该项之外的尾部全纯，故其主要部分不被更改。

\subsection{Mittag–Leffler 定理}
\label{b3:subsec:200403}

\begin{theorem}[Mittag--Leffler：整平面版本]\label{b3:thm:20-mittag-leffler}
任意无有限聚点的极点集及其有限主要部分，都由一个复平面亚纯函数实现。两个实现函数之差为整函数。
\end{theorem}
\begin{proof}
除去可能在原点的有限主要部分后，用上一小节的 Taylor 扣除构造并相加。局部正规收敛给亚纯性，极点邻域中的有限项分离给准确主要部分。两解的差在每个极点负幂全部抵消，因此所有极点可去，成为整函数。
\end{proof}
一般平面区域上的相应定理需要把多项式扣除改成适应补集的有理逼近；本节给出的完整构造范围是整平面，不把这一证明无说明地扩展到任意区域。

\subsection{与 Weierstrass 定理的对偶性}
\label{b3:subsec:200404}

Weierstrass 定理留下无零点的乘法自由度 \(e^g\)，Mittag--Leffler 定理留下整函数的加法自由度。两者还能组合给整插值：取简单零点为 \(a_n\) 的整函数 \(P\)，给主要部分 \(b_n/(P'(a_n)(z-a_n))\)，构造亚纯函数 \(G\)，则 \(F=PG\) 整且 \(F(a_n)=b_n\)。局部乘法直接核验值，任意增长的插值数据也不破坏离散点上的存在性。


主要部分构造也有误差与有效范围的区别。对每个 \(Q_n\) 减去有限 Taylor 多项式，减去的项不引入新极点，并且改变的是整函数自由度。若只减常数项，在近原点的紧集上尾部可能仍不收敛；所需次数应根据该项大小选择，而不是统一宣称一次修正足够。

\ProseParagraphBreak

整插值还可以指定有限阶导数。设在 \(a_n\) 希望 \(F\) 与给定多项式芽 \(J_n\) 的差至少有 \(m_n\) 阶零点。取 \(P\) 在 \(a_n\) 有恰好 \(m_n\) 阶零点，令 \(Q_n\) 是 \(J_n/P\) 的主要部分。由 Mittag--Leffler 构造 \(G\)，则 \(G-J_n/P\) 在 \(a_n\) 全纯，故 \(F=PG\) 满足 \(F-J_n\) 被 \((z-a_n)^{m_n}\) 整除。负幂计算准确保留了全部所需导数；不需要把每阶导数当作独立极点重新构造。

\ProseParagraphBreak

这种插值没有控制全平面增长。若要求在整数上取值 \(e^{n^2}\) 且整函数阶小于二，就会与最大模的下界矛盾；若指定的值只是有界，也不能自动保证存在有界整函数，因为 Liouville 使有界整函数只能恒常。存在性定理容许任意数据，正是因为它没有许诺这些额外估计。


\ProseParagraphBreak

下面把 Taylor 扣除步骤写成一个可检查的例子。在正整数 \(n\) 指定二阶主要部分 \(c_n/(z-n)^2\)，其中 \(c_n\) 可以任意迅速增长。对 \(|z|<n\)，有
\[
\frac{c_n}{(z-n)^2}=\frac{c_n}{n^2}\sum_{k\ge0}(k+1)(z/n)^k.
\]
取 \(R_n=n/2\)，选择整数 \(d_n\) 使
\[
\frac{|c_n|}{n^2}\sum_{k>d_n}(k+1)2^{-k}<2^{-n}.
\]
有限截断 \(T_n\) 为整多项式，故 \(\sum_n(c_n/(z-n)^2-T_n(z))\) 在任意避开正整数的紧集上正规收敛。对任意固定 \(R\)，仅有限个 \(R_n<R\)；剩余项由 \(2^{-n}\) 一致控制。由于没有一次负幂项，构造出的函数在每个极点的留数仍为零，不能因为修正次数高而凭空出现留数。

\ProseParagraphBreak

若 \(c_n=1\)，可以直接用 \(\sum_{n\ge1}(z-n)^{-2}\)，因为远处尾部已经是 \(O_K(n^{-2})\)；若 \(c_n=n^2\)，直接尾项趋于一，减去常数项后尾项仍有 \(O_K(n^{-1})\)，还须减一次项。具体地，
\[
\frac{n^2}{(z-n)^2}-1-\frac{2z}{n}=O_K(n^{-2}).
\]
由此可见，扣除次数应从具体尾项的展开决定。一般定理提供可用算法，特殊数据则可用更低次数的修正。

\ProseParagraphBreak

Mittag--Leffler 构造的唯一性比乘积构造更容易表达：具有相同极点主要部分的两个亚纯函数之差在每个极点可去，因而为整函数。反过来，加上任意整函数都不改变主要部分。因此“主要部分决定函数”只有在另有增长、周期性或规范化时才可能成立。对有限个极点，所有主要部分之和已经给出有理函数，但一般解仍是这个有理函数加任意整函数；极点数有限并不使函数必为有理。

\ProseParagraphBreak

插值点必须离散。例如若在 \(1/n\) 指定值零，同时在原点指定值一，任何全纯函数的连续性已产生矛盾；即使原点也指定零，恒等定理仍迫使函数恒为零，不能再在另一点任意指定非零值。有限阶导数插值也没有绕过这个限制。所谓“任意数据”准确地指在没有有限聚点的点列上，每一点只提出有限个导数条件。

% !TeX root = ../../Book3.tex
\section{Hadamard 分解}
\label{b3:sec:2005}

一般 Weierstrass 分解中的 \(g\) 可以任意整。若原函数有限阶，增长约束会把它压成多项式；这是 Hadamard 分解比单纯零点存在性更强的部分。

\subsection{有限阶整函数}
\label{b3:subsec:200501}

设整函数 \(f\) 阶为 \(\rho<\infty\)，非零零点 \(a_n\)，原点重数 \(m\)，令 \(p=\lfloor\rho\rfloor\)。零点计数 \(n(r)=O_\varepsilon(r^{\rho+\varepsilon})\)，故 \(\sum|a_n|^{-p-1}<\infty\)（取 \(\varepsilon<p+1-\rho\)）。因此 \(P=z^m\prod E_p(z/a_n)\) 存在，且 \(f/P=e^g\)，\(g\) 整。问题归结为证明 \(g\) 的导数高阶消失。

\subsection{典范乘积}
\label{b3:subsec:200502}

先说明所需的乘积估计。对每个小 \(\varepsilon>0\)，可以在每个充分大的区间 \([r,2r]\) 选一个半径 \(R\)，使
\[
\log|P(z)|\ge-C_\varepsilon R^{\rho+\varepsilon},\quad |z|=R.
\]
选半径时避开所有 \(|a_n|\le4r\) 的半径 \(r^{-\rho-2}\) 邻域；这些禁区总长为 \(O(r^{-2+\varepsilon/4})\)，小于 \(r\)。对内层零点，\(|1-z/a_n|\) 至少为所选间隔除 \(|a_n|\)，故对数下界为 \(-C\log r\)；有限的极小零点也满足同类下界。指数修正的 \(k\) 次项由分部求和控制，\(k\le p\le\rho\)，总量为 \(O(r^{\rho+\varepsilon/4}\log r)\)。外层 \(|a_n|>4r\) 用 \(|\log E_p(z/a_n)|\le2|z/a_n|^{p+1}\)，分部求和同样给 \(O(r^{\rho+\varepsilon/4})\)。最后用 \(\log r=O(r^{\varepsilon/2})\) 吸收对数。上界也由相同的分层估计得到。

关键不是逐点避开零点，而是选出整条无过近零点的圆周。否则 \(\log|P|\) 在零点处为负无穷，不能估计商的增长。

\subsection{Hadamard 分解定理}
\label{b3:subsec:200503}

\begin{theorem}[Hadamard 分解]\label{b3:thm:20-hadamard}
非恒零有限阶整函数 \(f\) 满足
\[
f(z)=z^m e^{Q(z)}\prod_nE_p(z/a_n),\quad p=\lfloor\rho(f)\rfloor,
\]
其中 \(Q\) 为次数不超过 \(\rho(f)\) 的多项式。
\end{theorem}
\begin{proof}
在上一小节选出的圆周，\(\operatorname{Re}g=\log|f|-\log|P|\le C_\varepsilon R^{\rho+\varepsilon}\)。调和最大值原理把界传到圆盘内部。
为把实部界变成模界，令 \(F=g-g(0)\)、\(A=\max_{|z|\le R}\operatorname{Re}F\)。若 \(A=0\)，最大值原理使 \(g\) 常数。若 \(A>0\)，函数 \(F/(2A-F)\) 在圆盘模不超过一，中心为零；Schwarz 引理给半径 \(R/2\) 上模不超过 \(1/2\)，解出 \(|F|\le2A\)。所以 \(\max_{|z|\le R/2}|g(z)|\le C_\varepsilon R^{\rho+\varepsilon}+|g(0)|\)。
Cauchy 估计对整数 \(k>\rho\) 给 \(|g^{(k)}(0)|\le k!C_\varepsilon R^{\rho+\varepsilon-k}\)。选 \(0<\varepsilon<k-\rho\)，沿选定半径趋无穷，得到导数零。于是 \(g=Q\) 为要求的多项式。
\end{proof}

\subsection{属与阶}
\label{b3:subsec:200504}

整函数的属通常取典范乘积属与 \(Q\) 次数的较大者。上述定理给属不超过有限阶；零点指数与指数多项式共同决定增长，不能只数零点。

作为应用，\(\sin\pi z\) 阶一、零点为全体整数。将正负因子配对，Hadamard 给
\[
\sin\pi z=e^{az+b}z\prod_{n\ge1}\left(1-\frac{z^2}{n^2}\right).
\]
两侧的奇性使 \(e^{az}\) 为偶函数，从 \(e^{az}=e^{-az}\) 对零点求导得到 \(a=0\)；再比较原点导数得 \(e^b=\pi\)。因此经典正弦乘积得到完整证明。


进一步展开乘积估计中的分部求和。若 \(n(t)\le Ct^a\)、\(a=\rho+\delta\)，则对 \(k\le p\le\rho\)，
\[
\sum_{1\le|a_n|\le4r}|a_n|^{-k}
\le (4r)^{-k}n(4r)+k\int_1^{4r}n(t)t^{-k-1}\,\mathrm{d}t
\le C_{k,\delta}r^{a-k}.
\]
乘上 \(r^k\) 后每阶修正均为 \(O(r^a)\)。当 \(p=0\) 没有修正项，内层只剩对数因子。对外层，取 \(a<p+1\)，同样得到
\[
\sum_{|a_n|>4r}|a_n|^{-p-1}\le
(p+1)\int_{4r}^{\infty}n(t)t^{-p-2}\,\mathrm{d}t
\le C r^{a-p-1}.
\]
这解释了为什么一定要先选 \(\delta<p+1-\rho\)：尾积分的指数只有在这个范围才可积。整数阶 \(\rho=p\) 仍有一整次的余量，而不是临界发散。

\ProseParagraphBreak

另一个可能的误用是先估计 \(|f/P|\)，再把分子分母的阶简单相减。整函数的阶并不这样运算；商在零点附近会遇到很小的分母，甚至两个同阶函数相减也可能造成完全抵消。本证明专门选择圆周上的乘积下界，再利用最大值原理，正是为了把分母风险变成可核验的估计。有限阶条件没有直接赋予商同样的阶。

\ProseParagraphBreak

实部转模的步骤是 \term[borel-caratheodory]{Borel--Carathéodory 估计}[Borel--Carathéodory estimate]的特例。更一般地，\(g\) 在闭半径 \(R\) 的圆盘全纯、\(A=\max\operatorname{Re}(g-g(0))\) 时，\(0<r<R\) 上有
\[
|g(z)-g(0)|\le\frac{2Ar}{R-r}.
\]
由同一分式 \(F/(2A-F)\) 与 Schwarz 估计 \(r/R\) 解出即可。此估计不是由 \(|\operatorname{Re}g|\) 控制全部函数；一个纯虚常数仍需由 \(g(0)\) 另外记录。 Hadamard 证明中取 \(r=R/2\)，使这一步的常数不随大半径改变。

\ProseParagraphBreak

有限阶且只有有限零点的整函数现在具有完全明确的形式 \(P e^Q\)。例如 \((z-1)^2e^{3z^2+z}\) 的多项式部分固定两个零点，指数二次项决定阶二、型三，一次项不改变阶和型。无零点整函数 \(e^{z^3}\) 的阶三全部来自指数因子。对 \(\sin\pi z\)，阶一由零点密度和双向指数增长同时确认；它没有附加线性指数因子，原因不是零点本身，而是奇偶对称。

\ProseParagraphBreak

在整数极点的对称例子中，配对能减少扣除。由正弦乘积的对数导数得
\[
\pi\cot\pi z=\frac1z+\sum_{n\ge1}\left(\frac1{z-n}+\frac1{z+n}\right)
=\frac1z+\sum_{n\ge1}\frac{2z}{z^2-n^2}.
\]
最后级数在避开整数的紧集正规收敛，因为尾部为 \(O_K(n^{-2})\)。把 \(z-n\) 与 \(z+n\) 两个极点单独按不对称顺序相加，却只得到 \(O(n^{-1})\) 的尾部。明确的配对是收敛证明的一部分，不能藏在省略号中。这个公式给所有整数留数一，同时把无穷多个局部极点资料汇总成一个周期函数。

\ProseParagraphBreak

最后检验两个同零点有限阶函数之间的关系。若 \(f_1,f_2\) 均非恒零、阶不超过 \(\rho\)，按同一零点列及同一 \(p=\lfloor\rho\rfloor\) 取乘积，分别有 \(f_j=P e^{Q_j}\)，\(\deg Q_j\le\rho\)。所以 \(f_1/f_2=e^{Q_1-Q_2}\)。这里并非先假设商有限阶，而是由两份分解共同推出。特别地，当 \(\rho<1\) 时商只能为常数；当 \(\rho=1\) 时仍可有 \(e^{az}\) 自由度。

\ProseParagraphBreak

阶零也包含非多项式整函数，不能把它误写成“阶零当且仅当多项式”。例如零点 \(a_n=e^{n^2}\) 的属零乘积满足
\[
\log M(r)\le\sum_{n\ge1}\log(1+r/e^{n^2})=O((\log r)^{3/2}),\quad r\ge e.
\]
取 \(N=\lfloor\sqrt{\log r}\rfloor\)，前 \(N+1\) 项各不超过 \(\log(1+r)\)，尾部由 \(r\sum_{n>N+1}e^{-n^2}\) 控制为有界量，故得到该估计。于是阶为零，但函数有无穷多个零点而不可能是多项式。Hadamard 在此只断言指数因子为常数；无穷乘积本身仍保留全部零点。有限阶与有限个零点是两项独立条件。


\ProseParagraphBreak

Hadamard 定理中的 \(p=\lfloor\rho\rfloor\) 是统一可用的选择，未承诺最小属。若零点临界幂和已经收敛，某些因子可降阶；降阶会改变指数多项式的低次项。例如正整数的属一乘积与一个事先配对的正负整数乘积拥有不同的指数修正，虽然最终零点可相同。比较两种表示时，应先把初等因子的指数差逐项展开，证明相关系数和收敛，再把它吸收进指数多项式。没有可和性就不能把一串指数形式地合并成一个有限系数。

\ProseParagraphBreak

正弦乘积还有一个实变量误差例子。对固定 \(|z|\le r\)，当 \(N>2r\) 时，尾因子 \(1-z^2/n^2\) 的对数绝对值总和不超过 \(2r^2\sum_{n>N}n^{-2}\le2r^2/N\)。于是截断乘积与完整乘积的尾部相对误差不超过 \(e^{2r^2/N}-1\)。离原点和整数零点很近时，函数自身可能小，相对全函数的计算误差还需考虑有限头部；在零点上，只能比较提出零因子后的函数。这给出了从解析恒等式到实际近似之间需要保留的条件。

\ProseParagraphBreak

由正弦乘积与倍角恒等式还可得到余弦乘积。将 \(\sin(2\pi z)/(2\sin\pi z)\) 按可去极点解释为 \(\cos\pi z\)，并在正规收敛的乘积中约去偶数零点因子，得
\[
\cos\pi z=\prod_{n\ge1}\left(1-\frac{4z^2}{(2n-1)^2}\right).
\]
约因子的合法性先在避开整数的紧集用非零尾部完成，再由全纯恒等定理扩展；不能在每个零点处直接把两个零乘积相除。原点值一确定常数因子，偶性排除一次指数自由度。该例同时展示零点删选、对称性和规范化怎样从一般分解确定一个具体函数。

% 本章小结（不计入编号）
\begin{chaptersummary}
对数可和性保证乘积尾部非零；初等因子使任意离散零点都能实现。扣除 Taylor 多项式把同一思想转成主要部分的加法构造。Hadamard 分解的额外责任是估计无零点商的实部，再用 Cauchy 估计证明其对数为多项式。正弦乘积展示了零点、增长与对称性共同确定特殊函数。
\end{chaptersummary}


\BookExercises

% 本题组自拟；各题新增训练动作见题意，未复制第三方原题。
\begin{exercise}\label{b3:ex:20-draft-01}
说明 \(\prod_{n\ge2}(1-z/n^2)\) 是否正规收敛，列出零点。
\end{exercise}
\begin{solution}
固定紧集半径 \(R\) 时 \(\sum R/n^2<\infty\)，正规乘积定理给整函数。有限零因子与非零尾部的分离说明零点恰为 \(n^2\)、\(n\ge2\)，均简单。
\end{solution}

\begin{exercise}\label{b3:ex:20-draft-02}
对零点 \(a_n=\sqrt n\) 求收敛指数与最小可用属。
\end{exercise}
\begin{solution}
幂和为 \(\sum n^{-s/2}\)，收敛当且仅当 \(s>2\)，所以指数二。\(p+1>2\) 才收敛，最小整数 \(p=2\)。
\end{solution}

\begin{exercise}\label{b3:ex:20-draft-03}
证明具有相同零点和重数的两个非恒零整函数之商为整指数。
\end{exercise}
\begin{solution}
商在共同零点处可去且无零点。对数导数整，在整平面有原函数 \(g\)；商乘 \(e^{-g}\) 导数零为非零常数，再把常数选一复对数并吸收进 \(g\)。
\end{solution}

\begin{exercise}\label{b3:ex:20-draft-04}
构造在每个整数 \(n\) 留数为一、没有其他极点的亚纯函数。
\end{exercise}
\begin{solution}
\(\pi\cot\pi z\) 在整数附近为 \(1/(z-n)+\)全纯项，所以符合。它由正弦、余弦整函数之商定义，其他点正弦非零，故无额外极点。
\end{solution}

\begin{exercise}\label{b3:ex:20-draft-05}
解释为何任何给定复数列 \(b_n\) 都能作为某整函数在正整数处的值。
\end{exercise}
\begin{solution}
正整数离散。取简单零点恰为正整数的整函数 \(P\)，用 Mittag--Leffler 构造主要部分 \(b_n/(P'(n)(z-n))\) 的 \(G\)。乘积 \(PG\) 的各极点可去，值为 \(b_n\)，没有增长限制。
\end{solution}

\begin{exercise}\label{b3:ex:20-draft-06}
证明无零点有限阶整函数形为 \(e^Q\)，且阶正数时必为整数。
\end{exercise}
\begin{solution}
Hadamard 乘积为空，且无原点零点，故为指数多项式。非恒定 \(Q\) 的次数 \(d\) 给最大模对数 \(|a_d|r^d+O(r^{d-1})\)，阶为整数 \(d\)。常数对应阶零。
\end{solution}

\begin{exercise}\label{b3:ex:20-draft-07}
由正弦乘积计算 \(\sum_{n\ge1}n^{-2}\)。
\end{exercise}
\begin{solution}
除以 \(\pi z\)，原点二阶系数由正弦级数为 \(-\pi^2/6\)。正规乘积展开的二阶系数为 \(-\sum n^{-2}\)，系数比较得 \(\pi^2/6\)。逐项取系数由局部一致收敛及 Cauchy 公式合法。
\end{solution}

\begin{exercise}\label{b3:ex:20-draft-08}
说明 \(e^{e^z}\) 有无零点，并判断其是否能写成指数多项式。
\end{exercise}
\begin{solution}
指数无零点，但函数无穷阶。若为 \(e^Q\)，其阶将为多项式次数，有限，矛盾。这显示 Hadamard 的有限阶条件不能删。
\end{solution}
